\documentclass[pdflatex,sn-mathphys-num]{sn-jnl}

\usepackage{graphicx}%
\usepackage{multirow}%
\usepackage{amsmath,amssymb,amsfonts}%
\usepackage{amsthm}%
\usepackage{mathrsfs}%
\usepackage[title]{appendix}%
\usepackage{xcolor}%
\usepackage{textcomp}%
\usepackage{manyfoot}%
\usepackage{booktabs}%
\usepackage{algorithm}%
\usepackage{algorithmicx}%
\usepackage{algpseudocode}%
\usepackage{listings}%
\usepackage{lmodern}%
\usepackage{siunitx}
\usepackage{microtype}
\usepackage{derivative}
\usepackage[version=4]{mhchem}
\usepackage{cleveref}
\usepackage[spanish]{babel}
\usepackage{braket}

\renewcommand{\arraystretch}{1.2}
\sisetup{group-digits=true,
  group-separator={\,},
  separate-uncertainty
}

\theoremstyle{thmstyleone}%
\newtheorem{theorem}{Teorema}%
\newtheorem{proposition}[theorem]{Proposición}%

\theoremstyle{thmstyletwo}%
\newtheorem{example}{Ejemplo}%
\newtheorem{remark}{Observación}%

\theoremstyle{thmstylethree}%
\newtheorem{definition}{Definición}%

\raggedbottom
\unnumbered% uncomment this for unnumbered level heads

\begin{document}

\title[Fundamentos de DFT y STM]{Características Matemáticas de la Teoría del
Funcional de la Densidad y Principios de la Microscopía de Efecto Túnel}

\author{\fnm{Julian L.} \sur{Avila-Martinez}}
\email{jlavilam@udistrital.edu.co}

\affil{\orgdiv{Programa Académico de Física}, \orgname{Universidad Distrital
Francisco José de Caldas}, \orgaddress{\city{Bogotá D.C.},
\country{Colombia}}}

\maketitle

\section{Características y Formulación Matemática de la Teoría del Funcional de
la Densidad (DFT)}

\subsection{El Espacio de Hilbert y el Problema de Muchos Cuerpos}

En la formulación estricta de la mecánica cuántica no relativista, el estado
puramente electrónico de un sistema de $N$ fermiones interactuantes se describe
mediante un vector de estado $\ket{\Psi}$ que pertenece al espacio de Hilbert
antisimetrizado $\mathcal{H}_N = \bigwedge_{i=1}^N L^2(\mathbb{R}^3) \otimes
\mathbb{C}^2$. La evolución estática del sistema está gobernada por el espectro
del hamiltoniano $H$, el cual, bajo la aproximación de Born-Oppenheimer,
se expresa a través del siguiente álgebra de operadores:
\begin{equation}
    H = T + V_{ee} + V_{\text{ext}}
\end{equation}
donde $T$ es el operador de energía cinética, $V_{ee}$ representa la interacción repulsiva de Coulomb
interelectrónica, y $V_{\text{ext}}$ es el potencial externo multiplicativo
(principalmente debido a la configuración de los núcleos).

La función de onda proyectada en la base de coordenadas $\Psi(\mathbf{r}_1, s_1, \dots, \mathbf{r}_N, s_N) = \braket{\mathbf{r}_1, s_1, \dots, \mathbf{r}_N, s_N | \Psi}$ depende explícitamente de $3N$ grados de libertad espaciales. Como se evidencia en sistemas físicos realistas, esta expansión dimensional prohíbe la resolución directa de la ecuación de autovalores $H\ket{\Psi} = E\ket{\Psi}$, dado que el costo computacional crece factorialmente.

\subsection{El Mapeo Funcional: Densidad y Vector de Estado} 

El aporte epistemológico y matemático fundamental de la Teoría del Funcional de la
Densidad radica en abandonar a $\ket{\Psi}$ como la variable primordial del
sistema. Los teoremas de Hohenberg-Kohn establecen una correspondencia biyectiva
estricta entre el potencial externo $v_{\text{ext}}(\mathbf{r})$ (salvo una constante
aditiva) y la densidad electrónica del estado fundamental $n(\mathbf{r})$.

Matemáticamente, la densidad se define como el valor esperado del operador
densidad de partículas:
\begin{equation}
    n(\mathbf{r}) = \braket{\Psi | n(\mathbf{r}) | \Psi} = N \sum_{s_1 \dots
    s_N} \int \lvert\Psi(\mathbf{r}, \mathbf{r}_2, \dots, \mathbf{r}_N)\rvert^2
    d^3r_2 \dots d^3r_N
\end{equation}

Dado que el potencial externo $v_{\text{ext}}(\mathbf{r})$ fija el hamiltoniano $H$, y este a su vez determina el estado fundamental $\ket{\Psi_0}$, existe un mapeo único $n(\mathbf{r}) \mapsto \ket{\Psi_0[n]}$. En consecuencia, toda propiedad observable del
estado fundamental es, por definición, un funcional matemático de la densidad
espacial $n(\mathbf{r})$. El funcional de energía total se reescribe
rigurosamente como
\begin{equation}
    E[n] = F_{HK}[n] + \int v_{\text{ext}}(\mathbf{r})n(\mathbf{r}) d^3r
\end{equation}
donde $F_{HK}[n] = \braket{\Psi[n] | T + V_{ee} | \Psi[n]}$ es un funcional universal independiente del sistema específico.

\subsection{Valor Computacional y Resolución Autoconsistente} 

El inmenso valor computacional de la DFT reside en que el principio variacional de Rayleigh-Ritz ya no se efectúa sobre el espacio $\mathcal{H}_N$ de $3N$ dimensiones, sino sobre el conjunto de funciones de densidad en $\mathbb{R}^3$. Esta compresión topológica hace que la tratabilidad de cristales y macro-moléculas sea viable.

Al desconocer la forma analítica exacta de los términos de intercambio y
correlación dentro del funcional, la resolución computacional requiere un
esquema iterativo no lineal conocido como campo autoconsistente (SCF). Conceptualmente, el procedimiento inicia con una densidad electrónica de prueba
empírica $n_{0}(\mathbf{r})$. Con ella, se evalúan los operadores efectivos del
sistema, permitiendo la diagonalización espectral del modelo de partículas
independientes. La solución de este sistema genera un nuevo conjunto de
autoestados que, al proyectarse, construyen una densidad actualizada $n_{1}(\mathbf{r})$.

Este ciclo algebraico se itera numéricamente ($\dots \to n_k \to n_{k+1} \to \dots$) hasta que la norma de la diferencia estocástica entre densidades sucesivas $\lVert n_{k+1}(\mathbf{r}) - n_k(\mathbf{r}) \rVert$ es menor a un umbral de tolerancia predefinido, garantizando que el sistema ha minimizado la hiper-superficie de energía del funcional.

\section{Principio de Funcionamiento del Microscopio de Efecto Túnel (STM)}

El microscopio de efecto túnel (STM) se aparta conceptualmente de las
herramientas ópticas convencionales; su funcionamiento no depende de la
refracción fotónica ni de interacciones macroscópicas, sino que explota
directamente la naturaleza probabilística de la mecánica cuántica a través del
efecto túnel.

En el diseño experimental, se aproxima una punta conductora (idealmente
fabricada en tungsteno y cortada de tal forma que termine geométricamente en un
único átomo) a la superficie de la muestra a distancias del orden de fracciones
de nanómetro, dentro de una cámara de ultra alto vacío. Desde la perspectiva
analítica clásica, el vacío entre la punta y el sustrato representa una barrera
de potencial infinita para los portadores de carga. Sin embargo, al resolver la
ecuación de Schrödinger para la función de onda de un electrón en la superficie
metálica, las soluciones permitidas presentan decaimientos exponenciales en la
región clásicamente prohibida (el vacío).

Al inducir una diferencia de potencial (voltaje de sesgo, $V$) entre la punta y
el sustrato, las funciones de onda evanescentes de ambos electrodos se solapan,
habilitando una probabilidad de transición distinta de cero. En la formulación
de teoría de perturbaciones dependiente del tiempo de Bardeen, la corriente de
tunelamiento resultante $I$ a bajas temperaturas se aproxima a
\begin{equation}
    I \propto \int_{0}^{eV} \rho_s(E) \rho_t(E - eV) \mathcal{T}(E, V, d) d E
\end{equation}
donde $\rho_s$ y $\rho_t$ son las densidades de estados del sustrato y la punta,
respectivamente, y $\mathcal{T}$ es el coeficiente de transmisión que decae
exponencialmente como $\sim \exp(-2\kappa d)$, siendo $d$ la distancia de
separación.

El rasgo físico fundamental del STM es que la imagen generada no representa una topografía física rígida; es una convolución espacial de la densidad de carga. Regiones interpretadas como ``luminosas'' u ``oscuras'' en la imagen obtenida
corresponden estrictamente a máximos y mínimos de la Densidad Local de Estados (LDOS). Esto permite elucidar tanto la conformación geométrica del retículo (o de las moléculas adsorbidas) como sus propiedades electrónicas superficiales con resolución sub-angstrom.

\end{document}
